% sec:geo-formulations — Lagrangian and Hamiltonian Formulations
\section{Lagrangian and Hamiltonian Formulations}
\label{sec:geo-formulations}

Every image produced by the ray tracer begins with the same computational
problem: given a photon's initial position $x^\mu_0$ and momentum
$p_{\mu,0}$, evolve the geodesic equation forward in affine parameter until
the ray hits something or escapes to infinity.  The formulation of that
equation matters.  \Cref{sec:dg-geodesic} derived it from two independent
starting points --- parallel transport and the variational principle --- and
showed that both the Lagrangian and Hamiltonian forms produce identical
trajectories.  The Lagrangian form requires all 40 Christoffel symbols per
evaluation; the Hamiltonian form needs only the metric, its inverse, and
first derivatives of the metric --- a difference that, for the Kerr
spacetime in Cartesian coordinates, separates a practical integrator from an
unmaintainable one.

\paragraph{The Lagrangian form and its cost.}

The geodesic Lagrangian (\cref{eq:geodesic-lagrangian}),
$L = \tfrac{1}{2}\,\metric\,\dot{x}^\mu\,\dot{x}^\nu$, yields the
geodesic equation through the Euler--Lagrange equations:
%
\begin{equation}
  \ddot{x}^\mu + \christoffel{\mu}{\alpha}{\beta}\,
  \dot{x}^\alpha\,\dot{x}^\beta = 0 \,.
  \label{eq:geo-lagrangian-eom}
\end{equation}
%
This reproduces \cref{eq:geodesic}, now viewed as the equation of motion that
the integrator must solve.  The system is second-order: four equations for
the four coordinate functions $x^\mu(\lambda)$.  Standard adaptive
integrators --- Runge--Kutta, Dormand--Prince, symplectic methods --- operate
on first-order systems, so the practical step is to introduce the velocity
$v^\mu = \dot{x}^\mu$ and rewrite as
%
\begin{equation}
  \dot{x}^\mu = v^\mu \,, \qquad
  \dot{v}^\mu = -\christoffel{\mu}{\alpha}{\beta}\,
  v^\alpha\,v^\beta \,.
  \label{eq:geo-lagrangian-first-order}
\end{equation}
%
This is a first-order system of eight equations with state vector
$(x^\mu, v^\mu)$.  The right-hand side of the velocity equation requires
all 40 independent Christoffel symbols at the current position.  Each
symbol (\cref{eq:christoffel-formula}) involves three partial derivatives
of the covariant metric and a contraction with the inverse metric.

For the Schwarzschild Kerr--Schild metric, where $f = 2M/r$ and
$l_\mu = (1, x/r, y/r, z/r)$, the required derivatives are manageable
analytically.  For the Kerr metric, where
$f = 2Mr^3/(r^4 + a^2 z^2)$ and the null covector $l_\mu$ involves the
oblate spheroidal relations of \cref{eq:kerr-oblate-spheroidal}, the
Christoffel expressions run to hundreds of terms in Cartesian
coordinates --- hand-coding them is an invitation to sign errors and
maintenance failures.
\warnnote{The first-order Lagrangian system
\cref{eq:geo-lagrangian-first-order} has a second drawback: the constraint
$\metric\,v^\mu v^\nu = \text{const}$ is not built into the equations.
Numerical error can change the norm of $v^\mu$, gradually turning a null
ray into a timelike or spacelike curve.}

\paragraph{The Hamiltonian form.}

The Legendre transform replaces the velocity $\dot{x}^\mu$ with the
conjugate momentum $p_\mu = \metric\,\dot{x}^\nu$
(\cref{eq:conjugate-momentum}).  The geodesic super-Hamiltonian
(\cref{eq:geodesic-hamiltonian}) is
%
\begin{equation}
  \mathcal{H} = \tfrac{1}{2}\,\inversemetric\,p_\mu\,p_\nu \,,
  \label{eq:geo-hamiltonian}
\end{equation}
%
and Hamilton's equations give eight first-order ODEs without auxiliary
variables:
%
\begin{align}
  \dot{x}^\mu &= \inversemetric\,p_\nu \,,
  \label{eq:geo-xdot} \\[4pt]
  \dot{p}_\mu &= -\tfrac{1}{2}\,(\pd{\mu}\,g^{\alpha\beta})\,
                 p_\alpha\,p_\beta \,.
  \label{eq:geo-pdot}
\end{align}
%
The position equation is algebraic: contract the inverse metric with the
momentum covector.  The momentum equation involves partial derivatives of the
inverse metric --- still a substantial computation, but one that can be
sidestepped entirely.

\paragraph{Derivatives of the covariant metric.}

An equivalent expression for $\dot{p}_\mu$ follows from the Euler--Lagrange
identity $\dot{p}_\mu = \partial L / \partial x^\mu$, derived in
\cref{sec:dg-geodesic}:
%
\begin{equation}
  \dot{p}_\mu = \tfrac{1}{2}\,(\pd{\mu}\,g_{\alpha\beta})\,
    \dot{x}^\alpha\,\dot{x}^\beta \,.
  \label{eq:geo-pdot-cov}
\end{equation}
%
This replaces derivatives of the \emph{inverse} metric with derivatives of
the \emph{covariant} metric.  The connection between the two forms runs
through the identity
$\pd{\mu}\,g^{\alpha\beta} =
-g^{\alpha\gamma}\,g^{\beta\delta}\,\pd{\mu}\,g_{\gamma\delta}$, obtained
by differentiating $g^{\alpha\gamma}\,g_{\gamma\beta} =
\delta^\alpha{}_\beta$; substituting into \cref{eq:geo-pdot} and replacing
$p_\alpha = g_{\alpha\gamma}\dot{x}^\gamma$ yields
\cref{eq:geo-pdot-cov}.  The velocity $\dot{x}^\alpha$ on the right-hand
side is not an independent variable --- it is computed from $p_\alpha$ via
\cref{eq:geo-xdot} at each step.
\physnote{Equations \cref{eq:geo-pdot,eq:geo-pdot-cov} are mathematically
equivalent.  The covariant form is computationally preferable because the
codebase stores the metric as a function of coordinates, making its
derivatives directly accessible via automatic differentiation.}

The Hamiltonian system in its final computational form ---
\cref{eq:geo-xdot,eq:geo-pdot-cov} --- requires three ingredients at each
integration step:
%
\begin{enumerate}
\item the inverse metric $\inversemetric$ (for the position equation),
\item the four partial-derivative matrices $\pd{\mu}\,g_{\alpha\beta}$
  (for the momentum equation),
\item the velocity $\dot{x}^\alpha = \inversemetric\,p_\nu$ (computed from
  the first ingredient and the current momentum).
\end{enumerate}
%
No Christoffel symbols appear.  The inverse metric is needed but never
differentiated.  For the Kerr--Schild spacetimes of
\cref{sec:schw-kerr-schild,sec:kerr-ks}, the inverse
$\eta^{\mu\nu} - f\,l^\mu l^\nu$ differs from the forward metric
$\eta_{\mu\nu} + f\,l_\mu l_\nu$ only by a sign on the $f\,l\,l$ term
(\cref{eq:ks-inverse,eq:kerr-ks-inverse}), so evaluating it costs nothing
extra.
\coderef{Spacetime.AutoDiff}{metricDerivativesAD}

The partial derivatives $\pd{\mu}\,g_{\alpha\beta}$ are computed by
automatic differentiation.  Each spacetime is implemented as a polymorphic
function with type signature \texttt{Floating a => [a] -> [a]},
taking four coordinates and returning ten independent metric components.
Passing an AD type instead of \texttt{Double} produces exact derivatives
$\pd{\mu}\,g_{\alpha\beta}$ without finite differences or hand-coded
expressions --- a strategy introduced in \cref{sec:dg-covariant} and
developed fully in \cref{ch:autodiff}.
\implnote{The function \texttt{metricDerivativesAD} computes all four
$\pd{\mu}\,g_{\alpha\beta}$ matrices in a single AD pass by extracting
columns of the Jacobian of the metric function.}

\paragraph{Three structural advantages.}

Beyond avoiding Christoffel symbols, the Hamiltonian formulation brings three
properties that the Lagrangian approach lacks:
%
\begin{enumerate}
\item \textbf{Natively first-order.}  The eight equations
  \cref{eq:geo-xdot,eq:geo-pdot-cov} are already in the form
  $\dot{y} = F(y)$ required by adaptive integrators, without
  introducing auxiliary velocity variables.

\item \textbf{Built-in error diagnostic.}  The super-Hamiltonian
  $\mathcal{H}$ is exactly conserved along any solution
  (\cref{eq:hamiltonian-constraint}): $\mathcal{H} = 0$ for null
  geodesics, $-\tfrac{1}{2}$ for timelike (with $\lambda$ equal to proper
  time and signature $({-}{+}{+}{+})$).  Any drift in
  $\mathcal{H}$ during numerical integration is pure discretisation
  error, providing a free accuracy check at every step
  (\cref{sec:geo-conservation}).

\item \textbf{Killing conservation.}  If the metric is independent of a
  coordinate $x^\alpha$, then $\pd{\alpha}\,g_{\mu\nu} = 0$ and
  \cref{eq:geo-pdot-cov} immediately gives $\dot{p}_\alpha = 0$: the
  corresponding momentum component is conserved.  This connects the
  Hamiltonian system directly to the Killing-vector analysis of the next
  section.
\end{enumerate}
%
The rest of this chapter works exclusively with
\cref{eq:geo-xdot,eq:geo-pdot-cov}.
